\documentclass[11pt]{article}

\usepackage[letterpaper,margin=0.9in]{geometry}
\usepackage[T1]{fontenc}
\usepackage[utf8]{inputenc}
\usepackage{amsmath,amssymb}
\usepackage{booktabs}
\usepackage{enumitem}
\usepackage{microtype}
\usepackage{setspace}
\usepackage{xcolor}
\usepackage{xurl}
\usepackage[colorlinks=true,linkcolor=blue!55!black,urlcolor=blue!55!black,
            citecolor=blue!55!black]{hyperref}

\setstretch{1.04}
\setlength{\parindent}{0pt}
\setlength{\parskip}{5pt}
\setlist{nosep,leftmargin=1.5em}
\newcommand{\code}[1]{\texttt{#1}}

\title{\vspace{-1.3cm}\textbf{Optimizing Defensive Assignments in\\
Co-Ed Slowpitch Softball}}
\author{David Russo}
\date{July 2026}

\begin{document}
\maketitle
\vspace{-0.8cm}

\begin{abstract}
This paper describes an optimization model for assigning available players to
defensive positions during a seven-inning co-ed slowpitch softball game. A
valid schedule must obey league rules about the number of fielders and the
number and placement of women, while also respecting each player's stated
position preferences. Among valid schedules, the model first seeks fair
playing time. It then tries to keep players at a small number of positions and
to give them continuous position assignments lasting two or more innings,
preferably three or more. The problem is formulated as a binary integer
constraint program and solved with Google's OR-Tools CP-SAT solver. A
two-stage objective separates playing-time fairness from positional
consistency. The resulting Python program is exposed through a mobile-friendly
Streamlit website so a manager can update last-minute availability and obtain
a new schedule in about three seconds.
\end{abstract}

\section{Problem Statement}

A slowpitch softball defense normally uses ten positions: pitcher (P),
catcher (C), first base (1B), second base (2B), third base (3B), shortstop
(SS), left field (LF), left center field (LC), right center field (RC), and
right field (RF). A roster may contain more than ten available players, so
some players must sit out in each inning. The manager must decide who plays,
who sits, and which position every fielder plays in each of seven innings.

The co-ed rules make this more than a simple rotation. A ten-player defense
requires at least four women. It must also include at least one woman in the
infield and at least one woman in the outfield. If only three women are
available, the team may use at most nine fielders, leaving RF empty. With only
eight total players, both RF and C are empty. Both eight- and nine-player
lineups require at least three women, including at least one in the infield and
one in the outfield. Pitcher and catcher count as infield positions.

Players also have unranked positional preferences. In this application a
preference is a hard eligibility rule: a player cannot be assigned to a
position that the player did not select. The one exception occurs when the
team fields eight or nine players. Because RF is then removed, a player who
selected RF is also considered eligible for RC. This exception gives a right
fielder a nearby position without weakening preferences in a normal
ten-player lineup.

A schedule should satisfy three practical goals in priority order:
\begin{enumerate}
    \item Every inning must be legal and every occupied position must be
    filled by an eligible player.
    \item Playing time should be as equal as the rules and preferences permit.
    \item Position assignments should be stable. A player should use few
    distinct positions and should remain at a newly assigned position for at
    least two consecutive innings, preferably three or more.
\end{enumerate}

These goals can conflict. For example, equal playing time might require a
specialist to enter for one inning, even though a longer assignment would be
more consistent. Therefore legality and preferences are modeled as hard
constraints, while fairness and consistency are modeled as ordered
optimization goals. The model optimizes one game at a time because
availability can change shortly before a game.

\section{Sets, Data, and Lineup Size}

Let $\mathcal{P}$ be the set of available players and let
$N=|\mathcal{P}|$. Let $\mathcal{W}\subseteq\mathcal{P}$ be the set of
available women. The seven innings are
$\mathcal{I}=\{1,\ldots,7\}$. The full position set is
\[
\mathcal{S}=\{\mathrm{P,C,1B,2B,3B,SS,LF,LC,RC,RF}\}.
\]
The infield and outfield subsets are
\[
\mathcal{S}_{IF}=\{\mathrm{P,C,1B,2B,3B,SS}\}, \qquad
\mathcal{S}_{OF}=\{\mathrm{LF,LC,RC,RF}\}.
\]

The lineup size $K$ and active position set $\mathcal{S}_K$ are computed
before the optimization model is built. This is useful because the league
rule determines $K$ directly from availability; the solver does not need to
choose it.

\begin{center}
\begin{tabular}{clc}
\toprule
$K$ & Active positions $\mathcal{S}_K$ & Minimum women $g_K$\\
\midrule
10 & all ten positions & 4\\
9  & all positions except RF & 3\\
8  & all positions except RF and C & 3\\
\bottomrule
\end{tabular}
\end{center}

The program selects $K=10$ if $N\geq 10$ and at least four women are
available. Otherwise it selects $K=9$ if $N\geq9$ and at least three women are
available. Otherwise it selects $K=8$ if $N=8$ and at least three women are
available. Fewer than eight players or fewer than three women produce an
infeasibility message before the solver is called.

For each player $p$, let $E_p\subseteq\mathcal{S}_K$ be that player's eligible
active positions. It is the intersection of the player's preferences and the
active position set. When $K\in\{8,9\}$ and RF was selected, RC is added to
$E_p$. A preflight check reports an error if no eligible player exists for a
required position or if no available woman is eligible for an active infield
or outfield position.

\section{Decision Variables}

The main decision variable is
\[
x_{ips} =
\begin{cases}
1, & \text{if player $p$ plays position $s$ in inning $i$},\\
0, & \text{otherwise.}
\end{cases}
\]
It is defined for $i\in\mathcal{I}$, $p\in\mathcal{P}$, and
$s\in E_p$. Variables for non-preferred positions are omitted completely,
which reduces model size and makes preference violations impossible.

Several auxiliary variables measure the quality of a schedule:
\begin{itemize}
    \item $t_p\in\{0,\ldots,7\}$ is the total innings played by player $p$.
    \item $y_{ps}$ is 1 if player $p$ uses position $s$ at least once.
    \item $u_p=\sum_s y_{ps}$ is the number of distinct positions used by
    player $p$.
    \item $e_p=\max(u_p-2,0)$ counts positions beyond the preferred soft limit
    of two.
    \item $a_{ips}$ is 1 when player $p$ starts a new consecutive stint at
    position $s$ in inning $i$.
    \item $q^{(1)}_{ips}$ is 1 when inning $i$ is an entire one-inning stint
    for player $p$ at position $s$.
    \item $q^{(2)}_{ips}$ is 1 when inning $i$ begins an exact two-inning stint
    for player $p$ at position $s$.
\end{itemize}

The model also uses $t^{\max}$ and $t^{\min}$ for the largest and smallest
individual playing-time totals. Their difference $R=t^{\max}-t^{\min}$ is the
playing-time range.

\section{Hard Constraints}

\subsection{Fill every active position}

Every active position must have exactly one player in every inning:
\[
\sum_{\substack{p\in\mathcal{P}\\s\in E_p}}x_{ips}=1
\qquad
\forall i\in\mathcal{I},\ s\in\mathcal{S}_K.
\]
This equation also implies that exactly $K$ players field in each inning.

\subsection{Use each player at most once per inning}

A person cannot occupy two positions in the same inning:
\[
\sum_{s\in E_p}x_{ips}\leq1
\qquad
\forall i\in\mathcal{I},\ p\in\mathcal{P}.
\]
Because $x_{ips}$ only exists for $s\in E_p$, this constraint works together
with the variable definition to enforce positional preferences.

\subsection{Meet the co-ed requirements}

Each inning must contain the required minimum number of women:
\[
\sum_{p\in\mathcal{W}}\sum_{s\in E_p}x_{ips}\geq g_K
\qquad \forall i\in\mathcal{I}.
\]
At least one woman must be in each defensive region:
\[
\sum_{p\in\mathcal{W}}
\sum_{s\in E_p\cap\mathcal{S}_{IF}}x_{ips}\geq1
\qquad \forall i\in\mathcal{I},
\]
\[
\sum_{p\in\mathcal{W}}
\sum_{s\in E_p\cap\mathcal{S}_{OF}}x_{ips}\geq1
\qquad \forall i\in\mathcal{I}.
\]
These are lower bounds rather than equalities, so a lineup may use more women
than the minimum.

\subsection{Define playing time and position use}

The innings played by each person are
\[
t_p=\sum_{i\in\mathcal{I}}\sum_{s\in E_p}x_{ips}.
\]
The largest and smallest values satisfy
$t^{\max}=\max_p t_p$ and $t^{\min}=\min_p t_p$. CP-SAT supports these
maximum and minimum relationships directly.

Position-use variables are linked to assignments by
\[
x_{ips}\leq y_{ps}\quad\forall i,p,s,
\qquad
y_{ps}\leq\sum_{i\in\mathcal{I}}x_{ips}\quad\forall p,s.
\]
Thus $y_{ps}$ equals 1 exactly when the position appears in that player's
schedule. The excess variable obeys $e_p\geq u_p-2$ and $e_p\geq0$.

\subsection{Identify consecutive position stints}

A stint is an uninterrupted run at the same position. Sitting out ends a
stint, even if the player later returns to the same position. For a start
variable after the first inning, the following linear inequalities represent
the logical statement ``assigned now and not assigned here previously'':
\[
a_{ips}\geq x_{ips}-x_{i-1,p,s},\quad
a_{ips}\leq x_{ips},\quad
a_{ips}\leq1-x_{i-1,p,s}.
\]
For the first inning, $a_{1ps}=x_{1ps}$.

For one-inning stints, define the boundary values $x_{0ps}=x_{8ps}=0$. Then
$q^{(1)}_{ips}$ represents the pattern $0,1,0$:
\[
\begin{aligned}
q^{(1)}_{ips}&\leq x_{ips},\\
q^{(1)}_{ips}&\leq1-x_{i-1,p,s},\\
q^{(1)}_{ips}&\leq1-x_{i+1,p,s},\\
q^{(1)}_{ips}&\geq x_{ips}-x_{i-1,p,s}-x_{i+1,p,s}.
\end{aligned}
\]
Similar ``if and only if'' inequalities identify an exact two-inning pattern
$0,1,1,0$. These variables do not prohibit short stints. Instead, they allow
the objective function to penalize them while preserving feasibility when
fairness or hard preferences make a short assignment unavoidable.

\section{Objective Function}

There is no single natural unit that compares fairness with consistency. A
one-inning stint, for example, cannot be directly compared with one inning of
playing-time imbalance. The implementation handles this by solving two
objectives in sequence.

\subsection{Stage 1: playing-time fairness}

There are $H=7K$ total fielding slots in a game. The perfectly equal share is
$H/N$ innings per player, which is usually fractional. CP-SAT uses integer
expressions, so the model defines the scaled absolute deviation
\[
D_p=\left|Nt_p-H\right|
    =N\left|t_p-\frac{H}{N}\right|.
\]
Multiplying by $N$ avoids fractions without changing which schedule is best.
Let $D^{UB}$ be a safe upper bound on $\sum_pD_p$. The first objective is
\[
\min\quad (D^{UB}+1)R+\sum_{p\in\mathcal{P}}D_p.
\]
The coefficient $D^{UB}+1$ makes the range more important than all possible
absolute deviations combined. Therefore the solver first minimizes the
difference between the most and least innings played. Among schedules with
that range, it minimizes total distance from the equal share.

\subsection{Stage 2: positional consistency}

After Stage 1, the best range and total deviation found are fixed as
constraints. Stage 2 cannot worsen them. It minimizes
\[
\begin{split}
1000\sum_{i,p,s}q^{(1)}_{ips}
&+200\sum_p e_p
+25\sum_{i,p,s}q^{(2)}_{ips}\\
&+5\sum_pu_p
+\sum_{i,p,s}a_{ips}.
\end{split}
\]
The weights express managerial preferences. A one-inning stint receives the
largest penalty. Using more than two distinct positions receives an additional
large penalty. Exact two-inning stints receive a smaller penalty, which means
three- or four-inning assignments are preferred when all else is equal. The
last two terms reduce the total number of positions and the number of times
position stints begin. Consequently, long stable assignments tend to win
remaining ties.

These are soft costs, not absolute rules. This distinction is important. With
specialists and strict gender requirements, eliminating every short stint can
be impossible without worsening playing-time fairness. The model returns the
best compromise it can find while never violating a hard rule.

\section{Solution Algorithm}

The model is implemented in Python with the OR-Tools CP-SAT solver. CP-SAT is
intended for optimization models made from integer and Boolean variables
\cite{ortools}. It combines constraint-programming ideas with Boolean
satisfiability methods. At a high level, it repeatedly propagates constraints
to rule out impossible assignments, branches on undecided variables, learns
from contradictions, and uses objective bounds to discard branches that
cannot improve the current solution. This is much more efficient than
explicitly listing every possible seven-inning schedule.

The application follows the procedure below.
\begin{enumerate}
    \item Read the roster, gender, availability, and preferences. Remove
    unavailable players and reject blank or duplicate names.
    \item Determine $K$, $\mathcal{S}_K$, $g_K$, and every eligibility set
    $E_p$. Run simple infeasibility checks before building the model.
    \item Create assignment variables and all hard constraints. Add the
    auxiliary fairness and consistency variables.
    \item Solve the Stage 1 fairness objective. With the current three-second
    overall limit, this phase receives up to about $0.75$ seconds.
    \item Fix the range and total deviation found in Stage 1. Pass its
    assignment as a solution hint, which gives Stage 2 a useful starting
    point.
    \item Use the remaining time to minimize the consistency penalty. The
    solver uses up to eight worker threads and a fixed random seed.
    \item If the status is \emph{OPTIMAL} or \emph{FEASIBLE}, translate the
    binary variables into a table. Otherwise, show an infeasibility or
    time-limit message.
\end{enumerate}

The time limit is a practical design choice. Assignment and scheduling
problems are combinatorial: adding players, innings, and eligible positions
can increase the search space very quickly. A game manager values a high
quality answer now more than a proof of optimality several minutes later.
Therefore a result marked \emph{FEASIBLE} is acceptable; it satisfies every
hard rule but the solver has not proved that no better objective value exists.
The website reports both solver status and elapsed time.

\section{Implementation, Output, and Limitations}

The website loads default names, genders, and preferences from
\code{roster\_positions.csv}. All players initially appear as available.
Users tap names to change availability and can edit preferences through
individual player cards. Because Streamlit keeps widget state separately for
each browser session, several low-volume users can run independent schedules
without a database.

The main output is a matrix with seven inning rows and the ten standard
position columns. Inactive positions contain a dash. A final \emph{Out} column
lists the available players sitting in that inning, separated by commas. A
second table reports total innings and positions by player, and the complete
matrix can be downloaded as a CSV file. Automated tests cover eight-, nine-,
and ten-player rules; the RF-to-RC exception; hard preferences; gender
placement; playing-time counts; position continuity; CSV defaults; and the
Streamlit interface.

The current model balances playing time only within one game. It does not
remember who sat more often in earlier games, and it assumes availability is
constant for all seven innings. Preferences are also unranked: every selected
position is equally acceptable. A future version could store season totals,
allow late arrival or early departure, distinguish ``preferred'' from
``allowed,'' or let a manager adjust objective weights. These additions would
require more input data and possibly persistent storage, but the same binary
assignment framework could support them.

\section{Conclusion}

The softball scheduling problem is a useful example of optimization because
it combines strict rules with human-centered preferences. Binary variables
represent who plays each position, linear constraints guarantee legal
lineups, and a two-stage objective gives clear priority to fairness before
positional stability. CP-SAT can search the resulting combinatorial space
quickly enough for last-minute roster changes. The deployed website turns the
mathematical model into a practical tool: a manager changes availability,
clicks one button, and receives a legal, fair, and reasonably consistent
seven-inning schedule.

\appendix
\section{Deployment to Streamlit Community Cloud}

The service was deployed to \emph{Streamlit Community Cloud}. Streamlit was
chosen because it hosts
small Python data applications without requiring a separately managed web
server, container cluster, or database. Community Cloud is available as a free
hosting service and connects directly to GitHub repositories
\cite{streamlitcloud}.

The deployment used the following steps:
\begin{enumerate}
    \item The local Python project was placed under Git version control. The
    root contains \code{app.py}, the optimizer package,
    \code{roster\_positions.csv}, \code{requirements.txt}, and the
    \code{.streamlit/config.toml} theme configuration.
    \item A GitHub repository was created at
    \url{https://github.com/davidluizrusso/softball-fielding}. The repository
    was made public so Streamlit could read it and the final app could be used
    by team members without requiring a login.
    \item Python dependencies were declared in \code{requirements.txt}:
    Streamlit for the interface, pandas for tabular data, and OR-Tools for
    CP-SAT. Community Cloud uses this file to recreate the Python environment
    \cite{streamlitdeps}.
    \item In the Community Cloud dashboard, \emph{Create app} was selected.
    The source was the \code{main} branch and the entry point was
    \code{app.py}. The equivalent GitHub file URL is
    \url{https://github.com/davidluizrusso/softball-fielding/blob/main/app.py}.
    Streamlit then installed the dependencies and launched the application
    \cite{streamlitdeploy}.
    \item The app's sharing setting was made public. Streamlit assigned a
    \code{streamlit.app} URL that can be opened in Chrome or Safari on desktop
    and mobile devices without an account.
\end{enumerate}

No secrets or database credentials are required. Roster changes made in the
page stay in that visitor's session, while the checked-in CSV provides the
default roster for a new session. Code updates follow a simple workflow:
modify and test locally, commit the changes, and run \code{git push}. Community
Cloud watches the GitHub repository and redeploys pushed changes automatically
\cite{streamlitmanage}. This approach meets the project's goals of low cost,
low administration, and a short path from local code to a shareable website.

\begingroup
\small
\setlength{\parskip}{2pt}
\begin{thebibliography}{9}
\bibitem{ortools}
Google for Developers. \href{https://developers.google.com/optimization/cp/cp_solver}
{\emph{CP-SAT Solver: OR-Tools}}. Accessed July 26, 2026.

\bibitem{streamlitcloud}
Streamlit. \href{https://docs.streamlit.io/deploy/streamlit-community-cloud}
{\emph{Streamlit Community Cloud}}. Accessed July 26, 2026.

\bibitem{streamlitdeps}
Streamlit.
\href{https://docs.streamlit.io/deploy/streamlit-community-cloud/deploy-your-app/app-dependencies}
{\emph{App Dependencies for Your Community Cloud App}}. Accessed July 26, 2026.

\bibitem{streamlitdeploy}
Streamlit.
\href{https://docs.streamlit.io/deploy/streamlit-community-cloud/deploy-your-app/deploy}
{\emph{Deploy Your App on Community Cloud}}. Accessed July 26, 2026.

\bibitem{streamlitmanage}
Streamlit. \href{https://docs.streamlit.io/deploy/streamlit-community-cloud/manage-your-app}
{\emph{Manage Your App}}. Accessed July 26, 2026.
\end{thebibliography}
\endgroup

\end{document}
